\begin{example}[Reachability]
In a directed graph with edge relation \(E\), the formula
\[
\begin{aligned}
  \operatorname{Reach}(a,b)\defeq
  \forall X\,\bigl[&
    \bigl(X(a)\land
      \forall u\,\forall v\,
      ((X(u)\land E(u,v))\to X(v))\bigr)\\
    &\to X(b)
  \bigr]
\end{aligned}
\]
holds under full semantics iff a finite directed path leads from
\(a\) to \(b\).  If a path exists, every edge-closed set containing \(a\)
contains its successive vertices.  Conversely, the actual set of vertices
reachable from \(a\) is edge-closed; full quantification permits it as a value
of \(X\).  The converse can fail in a Henkin structure if that reachable set
is not among the admissible unary relations.
\end{example}

\begin{warning}[Semantic parameter omitted]
Never write \(\Gamma\entails\varphi\) for a second-order argument before
declaring full or Henkin semantics.  The same sentence can have different
status under the two consequence relations.
\end{warning}

\begin{exercise}
\label{ex:second-order-logic-01}
Prove that the full second-order induction axiom is satisfied by the standard
natural numbers.  Your proof must quantify over an arbitrary subset of
\(\N\), not merely a definable subset.
\end{exercise}

\begin{exercise}
\label{ex:second-order-logic-02}
Explain precisely which step of the categoricity proof fails in a Henkin
structure whose unary-relation domain does not contain the range of \(f\).
\end{exercise}

\begin{exercise}
\label{ex:second-order-logic-03}
Use the compactness counterexample above to prove that no finitary sound and
strongly complete calculus exists for full second-order semantic consequence.
Distinguish this argument from the stronger non-enumerability theorem.
\end{exercise}

